\chapter{Conclusions} \label{ch:conclusions}

%This section should have a summary of the whole project.  The original aims and objectives, and whether they have been met, should be discussed. It should include a section that critiques and lists the limitations of your proposed solutions.  Future work should be described, and this should not be marginal or silly (e.g.\ add machine learning models).  It is always good to end on a positive note (i.e.\ `Final Remarks').

This study began with a single question: ``Why is instruction in Christian games sometimes considered bad?''. To answer this question, I first began by providing my own working definition of the ``Christian'' genre based on observations made by \textcite{gonzalez_born-again_2014, schut_swords_2026}. I described something (or someone) as Christian if it (or they) make a significant effort to follow the interpretation of any long-standing group or institution regarding the teachings of Jesus Christ. Then, I defined Christian games as ones with a Christian worldview, that is made evident through accessible signifiers. 

This allowed me to look at research related to this definition. However, I soon realised that only \textcite{schut_making_2013} had mentioned issues related to players' perceptions of instructive elements. Other works either sidestepped such issues entirely \parencite{gonzalez_born-again_2014}, or mostly neglected to consider the games' contents as a potential point of failure \parencite{schut_swords_2026, bernauer_playing_2012}. 

Therefore, I looked elsewhere for answers and found theories by \textcite{king_virtue_2017, repp_whats_2012}, each giving accounts of why heavy-handedness and heavy-handed instruction can be considered flaws, respectively; whereas the former proposes that didacticism might be an issue caused by perceived intellectual dishonesty, the latter blames problems concerning overtness on a restriction of aesthetic agency. I also proposed that the manner in which a game delivers instruction can also reduce engagement in the dimensions of the Player Involvement Model \parencite{calleja_in_game_2011}.

After confirming that the theories are compatible, I combined them into a single framework. I used it to produce analyses of two Christian games: \textit{Alum} \parencite{orsie_alum_2006} and \textit{Captain Bible} \parencite{engelbrite_captain_1994}. Through examples, I showed how the scripted narrative in \textit{Alum} \parencite{orsie_alum_2006} might lead some players to suspect intellectual shortcomings such as dogmatism and prejudice. However, I found it more difficult to show that \textit{Alum} restricts interpretative freedom. Although it might be a possible explanation, the relevant reviews I found could be explained in other ways. Meanwhile, in \textit{Captain Bible} I observed that its repetitive core gameplay loop, along with other formal elements, limits the affordances for ludic and affective involvement, making the levels feel vacant and dull. Therefore, players could disapprove of the game's attempts to edify. 

\section{Revisiting the Aims and Objectives}

Generally speaking, the framework proved quite practical in diagnosing the underlying issues in \textit{Alum} \parencite{orsie_alum_2006} and \textit{Captain Bible} \parencite{engelbrite_captain_1994}. Though I did not manage to make significant findings using the theory by \textcite{king_virtue_2017}, the warrant reduction assumption \parencite{repp_whats_2012} and Player Involvement Model \parencite{calleja_in_game_2011} proved to be quite useful. While I believe they are likely to apply to other digital games with similar issues, whether they are Christian or not, it would be unwise to make generalising statements given the small breadth of examples.

Regardless, for the examples in question, applying the framework helps clarify the issues indicated by players. Consequently, one might be able to suggest certain courses of action more confidently. For instance, given that \textit{Alum} \parencite{orsie_alum_2006} seems to appeal to more than just a Christian audience, Orsie could mitigate signs of cognitive vice by introducing a non-Christian playtester in his future projects. It is also possible that he could benefit from applying muse-based game design \parencite{khaled_muse-based_2012}. Not only would it allow him to get early input on his scripted narrative, but it would also require few resources, which might be appealing for an indie team like Crashable Studios.

\section{Critique and Limitations}


\section{Future Work}


\section{Final Remarks}

